\documentclass[11pt,a4paper]{article}
% Preambulo comun de los reportes de Tecnologias Celulares.
% Cada main.tex hace: \documentclass ... % Preambulo comun de los reportes de Tecnologias Celulares.
% Cada main.tex hace: \documentclass ... % Preambulo comun de los reportes de Tecnologias Celulares.
% Cada main.tex hace: \documentclass ... \input{../comun/preambulo.tex}
\usepackage[utf8]{inputenc}
\usepackage[T1]{fontenc}
\usepackage[spanish,mexico]{babel}
\usepackage{graphicx}
\usepackage[margin=2.5cm]{geometry}
\usepackage{amsmath}
\usepackage{amssymb}
\usepackage{booktabs}
\usepackage{tabularx}
\usepackage{array}
\usepackage{float}
\usepackage{setspace}
\usepackage{parskip}
\usepackage{fancyhdr}
\usepackage{textcomp}
\usepackage{xcolor}
\usepackage{caption}
\usepackage{fvextra}
\usepackage[hidelinks]{hyperref}

\onehalfspacing
\setlength{\parskip}{6pt}
\sloppy
\setlength{\emergencystretch}{3em}
\setlength{\headheight}{40pt}
\captionsetup{font=small,labelfont=bf}
\setlength{\tabcolsep}{6pt}

\pagestyle{fancy}
\fancyhf{}
\fancyhead[L]{\includegraphics[width=1cm]{/home/yovick/repos/UAZ-IEI/Assets/UIE.png}}
\fancyhead[R]{\small Tecnologías Celulares}
\fancyfoot[C]{\thepage}
\renewcommand{\headrulewidth}{0.4pt}

\DefineVerbatimEnvironment{codigo}{Verbatim}{%
  breaklines=true,breakanywhere=true,fontsize=\small,frame=single,
  framesep=3pt,rulecolor=\color{gray!50}}

% Figura centrada con pie: \figura[ancho]{archivo}{leyenda}
\newcommand{\figura}[3][0.85]{%
  \begin{center}
    \includegraphics[width=#1\linewidth]{#2}\\[2pt]
    {\small #3}
  \end{center}}

% Portada: \portada{numero}{titulo de la practica}
\newcommand{\portada}[2]{%
\begin{titlepage}
  \thispagestyle{empty}
  \begin{center}
    \vspace*{1cm}
    \includegraphics[width=4.5cm]{/home/yovick/repos/UAZ-IEI/Assets/uaz.png}
    \vspace{0.8cm}

    {\large Universidad Autónoma de Zacatecas}\\[4pt]
    {\large Unidad Académica de Ingeniería Eléctrica}\\[4pt]
    {\normalsize Ingeniería Electrónica Industrial con especialidad en Telecomunicaciones}

    \vspace{1.8cm}
    {\LARGE\bfseries Práctica #1}\\[10pt]
    {\Large\bfseries #2}
  \end{center}

  \vspace{2.2cm}
  \begin{center}
    \begin{tabular}{rl}
      \textbf{Materia:} & Tecnologías Celulares\\
      \textbf{Docente:} & Dr. Víktor Iván Rodríguez Abdalá\\
      \textbf{Alumno:}  & Arnold Yovick Ríos Zaldívar\\
      \textbf{Fecha:}   & Noviembre de 2026\\
    \end{tabular}
  \end{center}
  \vfill
\end{titlepage}}

\usepackage[utf8]{inputenc}
\usepackage[T1]{fontenc}
\usepackage[spanish,mexico]{babel}
\usepackage{graphicx}
\usepackage[margin=2.5cm]{geometry}
\usepackage{amsmath}
\usepackage{amssymb}
\usepackage{booktabs}
\usepackage{tabularx}
\usepackage{array}
\usepackage{float}
\usepackage{setspace}
\usepackage{parskip}
\usepackage{fancyhdr}
\usepackage{textcomp}
\usepackage{xcolor}
\usepackage{caption}
\usepackage{fvextra}
\usepackage[hidelinks]{hyperref}

\onehalfspacing
\setlength{\parskip}{6pt}
\sloppy
\setlength{\emergencystretch}{3em}
\setlength{\headheight}{40pt}
\captionsetup{font=small,labelfont=bf}
\setlength{\tabcolsep}{6pt}

\pagestyle{fancy}
\fancyhf{}
\fancyhead[L]{\includegraphics[width=1cm]{/home/yovick/repos/UAZ-IEI/Assets/UIE.png}}
\fancyhead[R]{\small Tecnologías Celulares}
\fancyfoot[C]{\thepage}
\renewcommand{\headrulewidth}{0.4pt}

\DefineVerbatimEnvironment{codigo}{Verbatim}{%
  breaklines=true,breakanywhere=true,fontsize=\small,frame=single,
  framesep=3pt,rulecolor=\color{gray!50}}

% Figura centrada con pie: \figura[ancho]{archivo}{leyenda}
\newcommand{\figura}[3][0.85]{%
  \begin{center}
    \includegraphics[width=#1\linewidth]{#2}\\[2pt]
    {\small #3}
  \end{center}}

% Portada: \portada{numero}{titulo de la practica}
\newcommand{\portada}[2]{%
\begin{titlepage}
  \thispagestyle{empty}
  \begin{center}
    \vspace*{1cm}
    \includegraphics[width=4.5cm]{/home/yovick/repos/UAZ-IEI/Assets/uaz.png}
    \vspace{0.8cm}

    {\large Universidad Autónoma de Zacatecas}\\[4pt]
    {\large Unidad Académica de Ingeniería Eléctrica}\\[4pt]
    {\normalsize Ingeniería Electrónica Industrial con especialidad en Telecomunicaciones}

    \vspace{1.8cm}
    {\LARGE\bfseries Práctica #1}\\[10pt]
    {\Large\bfseries #2}
  \end{center}

  \vspace{2.2cm}
  \begin{center}
    \begin{tabular}{rl}
      \textbf{Materia:} & Tecnologías Celulares\\
      \textbf{Docente:} & Dr. Víktor Iván Rodríguez Abdalá\\
      \textbf{Alumno:}  & Arnold Yovick Ríos Zaldívar\\
      \textbf{Fecha:}   & Noviembre de 2026\\
    \end{tabular}
  \end{center}
  \vfill
\end{titlepage}}

\usepackage[utf8]{inputenc}
\usepackage[T1]{fontenc}
\usepackage[spanish,mexico]{babel}
\usepackage{graphicx}
\usepackage[margin=2.5cm]{geometry}
\usepackage{amsmath}
\usepackage{amssymb}
\usepackage{booktabs}
\usepackage{tabularx}
\usepackage{array}
\usepackage{float}
\usepackage{setspace}
\usepackage{parskip}
\usepackage{fancyhdr}
\usepackage{textcomp}
\usepackage{xcolor}
\usepackage{caption}
\usepackage{fvextra}
\usepackage[hidelinks]{hyperref}

\onehalfspacing
\setlength{\parskip}{6pt}
\sloppy
\setlength{\emergencystretch}{3em}
\setlength{\headheight}{40pt}
\captionsetup{font=small,labelfont=bf}
\setlength{\tabcolsep}{6pt}

\pagestyle{fancy}
\fancyhf{}
\fancyhead[L]{\includegraphics[width=1cm]{/home/yovick/repos/UAZ-IEI/Assets/UIE.png}}
\fancyhead[R]{\small Tecnologías Celulares}
\fancyfoot[C]{\thepage}
\renewcommand{\headrulewidth}{0.4pt}

\DefineVerbatimEnvironment{codigo}{Verbatim}{%
  breaklines=true,breakanywhere=true,fontsize=\small,frame=single,
  framesep=3pt,rulecolor=\color{gray!50}}

% Figura centrada con pie: \figura[ancho]{archivo}{leyenda}
\newcommand{\figura}[3][0.85]{%
  \begin{center}
    \includegraphics[width=#1\linewidth]{#2}\\[2pt]
    {\small #3}
  \end{center}}

% Portada: \portada{numero}{titulo de la practica}
\newcommand{\portada}[2]{%
\begin{titlepage}
  \thispagestyle{empty}
  \begin{center}
    \vspace*{1cm}
    \includegraphics[width=4.5cm]{/home/yovick/repos/UAZ-IEI/Assets/uaz.png}
    \vspace{0.8cm}

    {\large Universidad Autónoma de Zacatecas}\\[4pt]
    {\large Unidad Académica de Ingeniería Eléctrica}\\[4pt]
    {\normalsize Ingeniería Electrónica Industrial con especialidad en Telecomunicaciones}

    \vspace{1.8cm}
    {\LARGE\bfseries Práctica #1}\\[10pt]
    {\Large\bfseries #2}
  \end{center}

  \vspace{2.2cm}
  \begin{center}
    \begin{tabular}{rl}
      \textbf{Materia:} & Tecnologías Celulares\\
      \textbf{Docente:} & Dr. Víktor Iván Rodríguez Abdalá\\
      \textbf{Alumno:}  & Arnold Yovick Ríos Zaldívar\\
      \textbf{Fecha:}   & Noviembre de 2026\\
    \end{tabular}
  \end{center}
  \vfill
\end{titlepage}}


\begin{document}

\portada{1}{Exploración de redes celulares utilizando RTL-SDR y MATLAB}

\textbf{Equipo:} RTL-SDR (\texttt{RTL2832U} + tuner \texttt{R820T}), antena telescópica

\textbf{Script:} \texttt{practica1.py} (Python con NumPy/SciPy; el manual pide MATLAB, pero
en este repositorio el procesamiento se hace en Python, que es software libre y produce
los mismos resultados numéricos)

\section{Objetivo}

Identificar y caracterizar señales de redes celulares presentes en el entorno mediante la
adquisición y el análisis de muestras IQ con un receptor RTL-SDR: localizar portadoras,
medir su frecuencia central y su ancho de banda, y comparar las distintas bandas celulares.

\section{Fundamento teórico}

Un receptor SDR entrega la envolvente compleja de la señal de RF ya trasladada a banda base:
\begin{equation}
  x[n] = I[n] + jQ[n],
\end{equation}
donde $I[n]$ es la proyección sobre $\cos(2\pi f_c t)$ y $Q[n]$ la proyección sobre
$-\sin(2\pi f_c t)$. La amplitud y la fase instantáneas son $|x[n]|$ y $\angle x[n]$.

El contenido espectral se obtiene con la DFT, que para $N$ muestras es
\begin{equation}
  X[k] = \sum_{n=0}^{N-1} x[n]\,e^{-j2\pi kn/N},
  \qquad
  f_k = f_c - \frac{F_s}{2} + k\frac{F_s}{N},
\end{equation}
con $F_s$ la frecuencia de muestreo. Para reducir la varianza del estimador se usa el
método de Welch: se promedian los periodogramas de segmentos solapados y se obtiene la
densidad espectral de potencia (PSD) en dBFS/Hz.

Una portadora LTE ocupa un ancho de banda nominal de 1.4, 3, 5, 10, 15 o 20 MHz, con
subportadoras de 15 kHz en el enlace descendente. Su PSD es aproximadamente plana dentro
del canal, con flancos de transición, y entre canales quedan bandas de guarda donde solo se
ve el piso de ruido. Esa diferencia de nivel es la que permite detectar cada portadora.

El RTL-SDR tiene tres limitaciones para este análisis: el ancho de banda instantáneo máximo
es de 3.2 MSPS (2.4 MSPS es el valor estable), el ADC es de 8 bits y la escala es relativa
al fondo de escala, no una potencia absoluta en dBm.

\section{Desarrollo experimental}

\subsection{Configuración del receptor}

\begin{tabularx}{\linewidth}{@{}*{2}{>{\raggedright\arraybackslash}X}@{}}
  \toprule
  Parámetro & Valor \\
  \midrule
  Frecuencia central de sintonía & 881.0 MHz (enlace descendente de 850 MHz) \\
  Frecuencia de muestreo & 20 MSPS en la captura sintética (2.4 MSPS en el RTL-SDR real) \\
  Muestras adquiridas & 262\,144 ($\approx$ 13.1 ms) \\
  Ganancia & AGC del tuner \\
  Antena & telescópica, banda ancha \\
  \bottomrule
\end{tabularx}

\subsection{Captura de muestras}

El equipo no tenía el RTL-SDR conectado durante la corrida, así que el script usa su modo
sintético: genera una captura IQ con el mismo formato que entregaría el receptor (piso de
ruido complejo, portadoras LTE con flancos de transición y dos espurios de banda angosta).
La cadena de análisis es idéntica a la que se aplica a una captura real; para usarla con
hardware basta reemplazar la fuente por el bloque \texttt{comm.SDRRTLReceiver}
(\texttt{pyrtlsdr} en Python) y fijar \texttt{--fs 2.4e6}. El reporte identifica en todo
momento qué resultados provienen de la captura sintética.

\subsection{Procesamiento}

\begin{enumerate}
  \item Se calcula la potencia promedio $P=\frac{1}{N}\sum |x[n]|^2$ en dBFS.
  \item Se estima la PSD con Welch (4096 muestras por segmento, 50\,\% de solapamiento).
  \item El piso de ruido se estima con el percentil 10 de la PSD; una portadora es una
        región contigua que supera el piso en 8 dB y mide al menos 0.4 MHz.
  \item Para cada región se integra la potencia y se recortan las colas al 0.5\,\% y
        99.5\,\% de la energía acumulada: el ancho entre esos dos puntos es el ancho de
        banda ocupado al 99\,\%.
  \item La frecuencia central se estima como el centroide de la potencia de la región.
  \item El espectrograma usa ventanas Hann de 1024 muestras con 75\,\% de solapamiento.
\end{enumerate}

\subsection{Exploración de bandas}

Se repite la cadena completa para las bandas de 700, 850, 1900, AWS 2100 y 2600 MHz,
cada una con su propia captura y su propio arreglo de portadoras.

\section{Código}

Fragmentos esenciales de \texttt{practica1.py}:

\begin{codigo}
def generar_captura(fc, fs, n, portadoras, piso=PISO_RUIDO_DBFS, semilla=7):
    rng = np.random.default_rng(semilla)
    x = (rng.standard_normal(n) + 1j * rng.standard_normal(n)) / np.sqrt(2)
    x *= 10 ** (piso / 20)
    for offset_mhz, ancho_mhz, potencia in portadoras:
        c = ruido_banda(n, fs, offset_mhz * 1e6, ancho_mhz * 1e6, rng)
        x += c * 10 ** (potencia / 20)
    return x

def detectar_portadoras(f, p_db, umbral_db=8.0, ancho_min=0.4e6):
    piso = float(np.percentile(p_db, 10))
    encima = p_db > piso + umbral_db
    ...
    acumulada = np.cumsum(pr) / np.sum(pr)
    i99 = np.searchsorted(acumulada, np.array([0.005, 0.995]))
    ancho = fr[i99[1]] - fr[i99[0]]
    centro = float(np.sum(fr * pr) / np.sum(pr))
\end{codigo}

\section{Resultados}

\subsection{Mediciones en la banda de 850 MHz}

\begin{tabularx}{\linewidth}{@{}*{2}{>{\raggedright\arraybackslash}X}@{}}
  \toprule
  Parámetro & Valor \\
  \midrule
  Potencia promedio de la captura & \textbf{$-$27.83 dBFS} \\
  Piso de ruido & \textbf{$-$145.0 dBFS/Hz} \\
  Relación pico a piso & 48.3 dB \\
  Portadora 1: frecuencia central & \textbf{877.006 MHz} \\
  Portadora 1: ancho de banda (99\,\%) & \textbf{10.03 MHz} \\
  Portadora 1: potencia integrada & $-$28.0 dBFS \\
  Portadora 2: frecuencia central & \textbf{887.011 MHz} \\
  Portadora 2: ancho de banda (99\,\%) & \textbf{5.11 MHz} \\
  Portadora 2: potencia integrada & $-$42.0 dBFS \\
  \bottomrule
\end{tabularx}

\subsection{Componentes I y Q}

\figura{figs/iq_temporal.png}{Componentes I y Q de la captura y su envolvente $|x[n]|$.}

Las componentes oscilan alrededor de cero con excursiones de $\pm0.07$ en la zona de las
portadoras. La envolvente no es constante porque la captura contiene varias señales
independientes y ruido; en un solo canal LTE la envolvente sería casi plana, salvo por los
picos propios de OFDM.

\subsection{Espectro y densidad espectral de potencia}

\figura{figs/espectro.png}{PSD de la captura. En magenta, las regiones detectadas automáticamente.}

Se distinguen dos mesetas planas separadas por una banda de guarda de 2.5 MHz. La meseta
ancha (10 MHz) corresponde a la portadora de mayor potencia y la angosta (5 MHz) a una
portadora vecina 10 dB más débil. El espurio de 883.4 MHz aparece como una línea delgada:
su ancho es menor que el mínimo de 0.4 MHz y por eso no se reporta como portadora.

\subsection{Espectrograma}

\figura{figs/espectrograma.png}{Espectrograma de los 13.1 ms de captura.}

Las dos portadoras se mantienen estables durante toda la ventana: no hay saltos de
frecuencia ni ráfagas, lo que es consistente con una emisión LTE continua. La banda de
guarda permanece en el piso de ruido durante todo el registro.

\subsection{Comparación entre bandas}

\figura[0.98]{figs/bandas.png}{Espectros de las cinco bandas exploradas. Cada sombra marca una portadora detectada.}

\begin{tabularx}{\linewidth}{@{}lcccc@{}}
  \toprule
  Banda & $f_c$ sint. & Portadoras detectadas & Anchuras & Potencias \\
        & (MHz) & (MHz) & (MHz) & (dBFS) \\
  \midrule
  700 MHz    & 740  & 733.01 / 744.00 & 5.11 / 5.11 & $-$44.0 / $-$38.0 \\
  850 MHz    & 881  & 877.01 / 887.01 & 10.04 / 5.11 & $-$28.0 / $-$42.0 \\
  1900 MHz   & 1950 & 1946.01 / 1956.01 & 10.04 / 5.11 & $-$40.0 / $-$34.0 \\
  AWS 2100   & 2140 & 2133.01 / 2143.00 & 5.11 / 10.03 & $-$39.0 / $-$31.0 \\
  2600 MHz   & 2650 & 2646.01 / 2656.01 & 10.04 / 5.11 & $-$46.0 / $-$37.0 \\
  \bottomrule
\end{tabularx}

Las anchuras medidas coinciden con las nominales dentro de 0.04 MHz, que es aproximadamente
la resolución de la PSD (4.9 kHz por bin) más el efecto del 99\,\% de potencia sobre los
flancos. La estimación de frecuencia central tiene un error menor a 15 kHz.

\section{Análisis de resultados}

\begin{itemize}
  \item \textbf{¿Qué bandas celulares fueron identificadas?} Las cinco bandas del manual:
        700, 850, 1900, AWS 2100 y 2600 MHz. En cada una se detectaron dos portadoras, con
        anchos de 5 y 10 MHz.
  \item \textbf{¿Qué diferencias existen entre las señales observadas?} Cambian la
        frecuencia central, el ancho de banda y el nivel de potencia. Las portadoras de
        10 MHz tienen una PSD más baja por hertzio que las de 5 MHz para la misma potencia
        total, porque reparten la energía en el doble de ancho de banda.
  \item \textbf{¿Cuál banda presentó mayor ocupación espectral?} En la captura de 850 MHz
        las portadoras ocupan 15.1 MHz de los 20 MHz observados (75\,\%). La misma
        fracción aparece en 1900, AWS y 2600 (15.1 a 15.2 MHz); solo 700 MHz queda en
        10.2 MHz (51\,\%) porque tiene dos canales de 5 MHz.
  \item \textbf{¿Se observaron otros servicios además de LTE?} Sí: los espurios de banda
        angosta en 883.4 y 877.1 MHz (tonos de 1.7 MHz y $-$3.9 MHz respecto de la
        sintonía) no tienen la forma plana de un canal LTE y se descartan por ancho.
  \item \textbf{Limitación del modo sintético.} Las cifras anteriores describen la captura
        generada, no el entorno radioeléctrico real de Zacatecas. Con el dongle conectado
        los anchos medidos dependerían de la ocupación real y de la atenuación del sitio;
        el algoritmo de detección es el mismo.
\end{itemize}

\section{Cuestionario}

\begin{enumerate}
  \item \textbf{¿Qué es una red celular?} Un sistema de radiocomunicación que divide el
        territorio en celdas, cada una servida por una estación base, y reutiliza las
        frecuencias a cierta distancia para dar servicio móvil a muchos usuarios con un
        espectro limitado.
  \item \textbf{¿Qué diferencia existe entre 2G, 3G y 4G?} 2G (GSM) usa TDMA/FDMA y voz
        digital; 3G (UMTS) usa CDMA de banda ancha y añade datos a mayor tasa; 4G (LTE)
        usa OFDM/OFDMA con MIMO y conmutación de paquetes, con tasas de decenas a cientos
        de Mbps y latencias menores.
  \item \textbf{¿Por qué se utilizan bandas licenciadas?} Porque el Estado asigna bloques
        exclusivos a cada operador; así se evita la interferencia entre redes y se puede
        planificar la cobertura y la calidad del servicio. Las bandas libres (ISM) no
        ofrecen esa garantía.
  \item \textbf{¿Qué información puede obtenerse a partir del espectro?} La frecuencia
        central y el ancho de banda de cada emisión, su potencia relativa, la separación
        entre canales, la presencia de guardas, interferencias y espurios, y la ocupación
        de la banda en el tiempo.
  \item \textbf{¿Qué limitaciones tiene el RTL-SDR para analizar señales LTE?} Su ancho de
        banda instantáneo (2.4--3.2 MHz) es menor que un canal LTE de 10 o 20 MHz, su ADC
        de 8 bits da poco rango dinámico y su escala no está calibrada en dBm. Además, el
        oscilador sin TCXO introduce decenas de ppm de error en frecuencia.
\end{enumerate}

\section{Conclusiones}

Se implementó la cadena completa de exploración: captura IQ, estimación de potencia, PSD
de Welch, espectrograma y detección automática de portadoras. En la banda de 850 MHz se
identificaron dos portadoras en 877.006 y 887.011 MHz, con anchos de 10.03 y 5.11 MHz y
potencias de $-$28.0 y $-$42.0 dBFS sobre un piso de $-$145.0 dBFS/Hz. El mismo
procedimiento aplicado a 700, 1900, AWS y 2600 MHz detectó dos portadoras por banda, con anchos de 5 y 10 MHz (dos de 5 MHz en 700 MHz). La captura es sintética porque no había receptor conectado;
el script conserva la interfaz para sustituirla por muestras reales sin cambiar el análisis.

\end{document}
